%! Semi-log Plots — Unit 2, Lesson 2.15 — Slides (Beamer)
\documentclass[aspectratio=169, 11pt]{beamer}
\usepackage{precalculus-beamer}   % provides \navyheader{} and \sectionlabel[color]{}

\begin{document}

% TITLE SLIDE
{
\setbeamercolor{background canvas}{bg=navy}
\begin{frame}[plain]
  \vspace{2.8cm}\hspace{0.55cm}%
  \begin{minipage}{0.9\textwidth}
    {\color{goldacc}\bfseries\small LESSON 2.15}\\[0.5cm]
    {\color{white}\bfseries\LARGE Semi-log Plots}\\[1.2cm]
    {\color{skymid}\small Precalculus~$\cdot$~Unit 2: Exponential and Logarithmic Functions}
  \end{minipage}
\end{frame}
}

% SLIDE 1 — HOOK
\begin{frame}[t, plain]
  \navyheader{The Scientist's Dilemma}
  \vspace{0.3cm}
  \begin{block}{Consider}
    A scientist plots bacterial population data on two graphs:
    \begin{itemize}
      \item \textbf{Regular graph:} the curve bends sharply upward.
      \item \textbf{Semi-log graph:} the same data forms a perfectly straight line.
    \end{itemize}
    \vspace{0.2cm}
    \textbf{What does the straight line on the semi-log graph tell the scientist?}
  \end{block}
  \vspace{0.4cm}
  \textit{Discuss with a partner for 1 minute, then share.}
\end{frame}

% SLIDE 2 — WHAT IS A SEMI-LOG PLOT?
\begin{frame}[t, plain]
  \navyheader{What Is a Semi-log Plot? (EK 2.15.A.1)}
  \vspace{0.2cm}
  \begin{block}{Definition}
    A \textbf{semi-log plot} has one axis (usually the $y$-axis) that is
    \textbf{logarithmically scaled}: equal spacing represents equal \emph{ratios}
    (powers of the base), not equal differences.
  \end{block}
  \vspace{0.3cm}
  \begin{itemize}
    \setlength{\itemsep}{8pt}
    \item On a semi-log plot, data that follows $y = ab^x$ \textbf{appears linear}.
    \item Non-exponential data \textbf{curves} on a semi-log plot.
    \item \textbf{Advantage (EK 2.15.A.2):} No constant needs to be added to $y$-values
          to check for exponential behavior.
  \end{itemize}
\end{frame}

% SLIDE 3 — LINEARIZING y = ab^x
\begin{frame}[t, plain]
  \navyheader{Linearizing $y = ab^x$ (EK 2.15.B.2)}
  \vspace{0.2cm}
  Apply $\log_{10}$ to both sides of $y = ab^x$:
  \[
    \log_{10}(y) = \log_{10}(a) + x\log_{10}(b)
  \]
  \begin{block}{The Linearized Model}
    \[
      \log_{10}(y) = \underbrace{(\log_{10} b)}_{\text{slope}}\cdot x
                   + \underbrace{\log_{10} a}_{\text{intercept}}
    \]
  \end{block}
  \vspace{0.2cm}
  \textbf{Example:} $y = 5\cdot3^x$
  \begin{itemize}
    \setlength{\itemsep}{4pt}
    \item Slope $= \log_{10}(3) \approx 0.477$
    \item Intercept $= \log_{10}(5) \approx 0.699$
    \item Linearized form: $\log_{10}(y) \approx 0.477x + 0.699$
  \end{itemize}
\end{frame}

% SLIDE 4 — RECOVERING THE EXPONENTIAL MODEL
\begin{frame}[t, plain]
  \navyheader{Recovering the Exponential Model (EK 2.15.B.1)}
  \vspace{0.2cm}
  Given slope $m$ and $y$-intercept $c$ from a semi-log plot:
  \[
    b = 10^m \qquad \text{and} \qquad a = 10^c
  \]
  \begin{block}{Example}
    A semi-log plot of $(x,\,\log_{10}(y))$ has slope $0.602$ and $y$-intercept $0.699$.
    \begin{itemize}
      \setlength{\itemsep}{4pt}
      \item $b = 10^{0.602} \approx 4$
      \item $a = 10^{0.699} \approx 5$
      \item Exponential model: $y = 5\cdot4^x$
    \end{itemize}
  \end{block}
\end{frame}

% SLIDE 5 — ADVANTAGE OVER ADDING CONSTANTS
\begin{frame}[t, plain]
  \navyheader{Advantage of Semi-log Plots (EK 2.15.A.2)}
  \vspace{0.3cm}
  \begin{columns}
    \column{0.48\textwidth}
    \begin{block}{Lesson 2.6 Approach}
      \begin{itemize}
        \setlength{\itemsep}{6pt}
        \item Add a constant $k$ to all $y$-values.
        \item Check if the adjusted values show proportional growth.
        \item Requires guessing a good value for $k$.
      \end{itemize}
    \end{block}
    \column{0.48\textwidth}
    \begin{block}{Semi-log Approach (Today)}
      \begin{itemize}
        \setlength{\itemsep}{6pt}
        \item Compute $\log_{10}(y)$ for each data point.
        \item Check if $\log_{10}(y)$ vs.\ $x$ is linear.
        \item \textbf{No constant required at all.}
      \end{itemize}
    \end{block}
  \end{columns}
  \vspace{0.3cm}
  \centering
  \textit{If the semi-log data is linear, an exponential model is appropriate --- guaranteed.}
\end{frame}

% SLIDE 6 — CLASS PRACTICE
\begin{frame}[t, plain]
  \navyheader{Class Practice}
  \vspace{0.2cm}
  \textbf{Given:} $x = 0,1,2,3$ and $y = 4, 12, 36, 108$.
  \vspace{0.2cm}
  \begin{enumerate}
    \setlength{\itemsep}{8pt}
    \item Compute $\log_{10}(y)$ for each $x$.
    \item What is the slope of $\log_{10}(y)$ vs.\ $x$?
    \item What is the $y$-intercept?
    \item Find $b$ and $a$, then write the exponential model $y = a\cdot b^x$.
    \item Verify your model with $y(1)$.
  \end{enumerate}
  \vspace{0.3cm}
  \begin{block}{Key Formulas}
    Slope $= \log_{10}(b)$ \quad $\Rightarrow$ \quad $b = 10^{\text{slope}}$
    \qquad
    Intercept $= \log_{10}(a)$ \quad $\Rightarrow$ \quad $a = 10^{\text{intercept}}$
  \end{block}
\end{frame}

\end{document}
